% ============================================================
% РАЗДЕЛ 5. Оценка эффективности и надёжности
% ============================================================
\section{ОЦЕНКА ЭФФЕКТИВНОСТИ И НАДЁЖНОСТИ ПРЕДЛОЖЕННОГО ПОДХОДА}

В данном разделе описывается методика экспериментальной оценки предложенной архитектуры, формулируются метрики, строится расчётная модель задержки, анализируются ожидаемые результаты сравнения и рассматриваются ограничения подхода.

\subsection{Методика экспериментов}

\subsubsection*{Тестовый стенд}

Для сравнительной оценки предлагается следующая конфигурация тестового стенда. Кластер состоит из 6 worker-нод на базе x86\_64-серверов с процессорами Intel Xeon Silver (16 физических ядер, 32 потока), 64~ГБ ОЗУ и локальным NVMe-диском Samsung 980 Pro ёмкостью 2~ТБ (последовательная скорость чтения до 7000~МБ/с). Сеть между нодами --- 10~Гбит/с Ethernet с коммутатором Top-of-Rack, RTT между нодами не превышает 0.5~мс.

Redis развёртывается в конфигурации Redis Sentinel: 1 master, 2 replica, автоматический failover при недоступности master. Размер чанка зафиксирован на уровне 16~МБ как компромисс между накладными расходами на метаданные (меньший чанк --- больше записей в Redis) и гранулярностью параллельной загрузки (более крупный чанк снижает степень P2P-параллелизма).

Набор тестовых моделей включает: Llama~3~8B ($\approx 15$~ГБ, $\approx 960$ чанков), Mistral~7B ($\approx 14$~ГБ, $\approx 896$ чанков) и Llama~3~70B ($\approx 140$~ГБ, $\approx 8960$ чанков). Для воспроизводимости результатов в отсутствие GPU runtime-контейнер заменяется синтетическим читателем, который последовательно читает все файлы модели через FUSE-точку монтирования и фиксирует время до первого байта и общее время чтения.

\subsubsection*{Процедура эксперимента}

Каждый сценарий тестируется в двух режимах: \textit{cold cluster} (дисковые кэши очищены, Redis пуст) и \textit{warm cluster} (модель загружена хотя бы на одной ноде). В режиме warm cluster дополнительно разделяются сценарии: warm on same node (модель на ноде, где стартует под) и warm on peer node (модель на другой ноде кластера). Для burst-сценария одновременно запускаются 10 подов на разных нодах, и измеряется суммарный внешний трафик и распределение $T_{\mathrm{model\_ready}}$ по репликам.

\subsubsection*{Сравниваемые сценарии}

Сравниваются три подхода:

\begin{enumerate}[leftmargin=*]
  \item \textbf{Baseline} (\texttt{storage-initializer + emptyDir}): каждый под независимо загружает всю модель из HuggingFace Hub через публичный интернет-канал. Представляет текущую практику в кластерах без выделенного кэша.

  \item \textbf{NFS}: общая директория монтируется через PVC/NFS. Первый под загружает модель из репозитория в NFS, остальные читают напрямую из общего тома. Представляет промежуточный подход, описанный в разделе~3.

  \item \textbf{P2P/FUSE} (предложенный подход): каждый под монтирует виртуальную файловую систему через \texttt{storage-mounter}; FUSE-операции транслируются в chunk-запросы к \texttt{storage-agent}. Агент обслуживает запросы из local cache, затем из peer cache (TCP к соседним агентам), и только при отсутствии чанков в кластере --- из HuggingFace (remote fallback).
\end{enumerate}

Каждый эксперимент повторяется не менее 5 раз; фиксируются медиана и 99-й перцентиль (p99) каждой метрики. Между экспериментами кэши очищаются или намеренно прогреваются в зависимости от измеряемого сценария. Результаты записываются средствами Prometheus и Grafana; временна́я метка создания пода берётся из Kubernetes Event API.

\subsection{Метрики оценки}

Для каждого сценария измеряются следующие метрики.

\textbf{Время до первого байта} $T_{\mathrm{first\_byte}}$ --- интервал от момента создания пода (pod creation timestamp) до момента, когда runtime-контейнер успешно считал первый байт файла весов. Эта метрика критична для реализации lazy loading, при котором инференс может начаться до полной загрузки модели.

\textbf{Время готовности модели} $T_{\mathrm{model\_ready}}$ --- полное время от создания пода до момента, когда синтетический читатель завершил последовательное чтение всех файлов модели. Для реального ML-runtime эквивалентно времени загрузки весов в GPU-память.

\textbf{Cache hit rate} $H_l$ определяется как доля запросов к \texttt{storage-agent}, обслуженных из локального дискового кэша:
\[
  H_l = \frac{N_{\mathrm{local\_hit}}}{N_{\mathrm{local\_hit}} + N_{\mathrm{peer\_hit}} + N_{\mathrm{remote}}}
\]

\textbf{Peer hit rate} $H_p$ --- доля запросов, удовлетворённых с соседних агентов по TCP:
\[
  H_p = \frac{N_{\mathrm{peer\_hit}}}{N_{\mathrm{local\_hit}} + N_{\mathrm{peer\_hit}} + N_{\mathrm{remote}}}
\]

\textbf{Remote fallback rate} $R$ = $1 - H_l - H_p$ --- доля запросов, потребовавших обращения к внешнему репозиторию.

\textbf{Внешний сетевой трафик} $V_{\mathrm{ext}}$ --- суммарный объём данных, скачанный из HuggingFace при масштабировании до $N$ реплик. Для baseline: $V_{\mathrm{ext}} = N \times S_{\mathrm{model}}$; для P2P/FUSE в идеальном случае: $V_{\mathrm{ext}} = S_{\mathrm{model}}$ (однократная загрузка).

\textbf{Redis lookup latency} --- время выполнения команд \texttt{SMEMBERS chunk:peers:\{hash\}} (поиск пиров для чанка) и \texttt{HGET agent:\{id\}} (получение адреса агента), измеренное как p50 и p99. Ожидаемые значения: p50 $<$ 1~мс, p99 $<$ 5~мс при развёртывании Redis в той же сети кластера~\cite{redis_sentinel_docs}.

\textbf{CPU overhead FUSE} --- дополнительная нагрузка на CPU worker-ноды, вносимая FUSE-демоном \texttt{storage-mounter}, выраженная в процентах от доступных ядер при последовательном чтении 16-МБ чанков.

\textbf{Задержка репликации чанка} --- время от момента записи чанка первым агентом до его появления в реестре Redis для других агентов. Включает время записи на NVMe, время команды \texttt{SADD} к Redis и RTT до Redis-сервера. Ожидаемое суммарное значение не превышает $T_l + \mathrm{RTT_{Redis}} \approx 2.3 + 1 \approx 3.3$~мс.

\subsection{Расчётная модель задержки}

\subsubsection*{Трёхуровневая модель AMAT}

Среднее время доступа к данным (AMAT, Average Memory Access Time) для трёхуровневой иерархии кэша~\cite{hennessy2011ca} определяется формулой:
\begin{equation}
  T_{\mathrm{AMAT}} = H_l \cdot T_l + (1 - H_l)\,[\,H_p \cdot T_p + (1 - H_p) \cdot T_r\,]
  \label{eq:amat-evaluation}
\end{equation}
где $T_l$ --- время чтения из локального дискового кэша, $T_p$ --- время получения чанка от соседнего агента по TCP, $T_r$ --- время загрузки чанка из внешнего репозитория, $H_l$ и $H_p$ --- доли попаданий на каждом уровне.

\subsubsection*{Характерные значения компонентов задержки}

Для NVMe-диска при последовательном чтении чанка размером 16~МБ:
\[
  T_l = \frac{16 \times 10^6 \text{ байт}}{7000 \times 10^6 \text{ байт/с}} \approx 2.3~\text{мс}
\]

Для P2P-передачи чанка по сети 10~Гбит/с с RTT 0.5~мс:
\[
  T_p = \frac{16 \times 10^6 \times 8~\text{бит}}{10^{10}~\text{бит/с}} + \mathrm{RTT} \approx 12.8~\text{мс} + 0.5~\text{мс} \approx 13.3~\text{мс}
\]

Для загрузки чанка из HuggingFace Hub через типичный канал 200~МБ/с (с учётом перегрузки сервера):
\[
  T_r = \frac{16 \times 10^6~\text{байт}}{200 \times 10^6~\text{байт/с}} + T_{\mathrm{internet-RTT}} \approx 80~\text{мс} + 50~\text{мс} \approx 130~\text{мс}
\]
При перегруженном репозитории или медленном внешнем канале $T_r$ может достигать 500--5000~мс.

\subsubsection*{Расчёт AMAT при различных сценариях}

Подставляя характерные значения $T_l \approx 2.3$~мс, $T_p \approx 13.3$~мс, $T_r \approx 130$~мс в формулу~(\ref{eq:amat-evaluation}):

\begin{table}[htbp]
\centering
\caption{Расчётные значения AMAT при разных уровнях hit rate}
\label{tab:amat_calc}
\begin{tabular}{cccp{4cm}r}
\toprule
$H_l$ & $H_p$ & $R$ & Описание сценария & AMAT (мс) \\
\midrule
0.00 & 0.00 & 1.00 & Первый запуск, нет кэша & $\approx 130$ \\
0.00 & 0.80 & 0.20 & Нет локального, есть соседи & $\approx 37$ \\
0.90 & 0.80 & 0.02 & После прогрева кластера & $\approx 4.3$ \\
0.99 & 0.80 & 0.002 & Горячий кэш, популярная модель & $\approx 2.5$ \\
\bottomrule
\end{tabular}
\end{table}

Таблица~\ref{tab:amat_calc} демонстрирует, что после прогрева кластера (hit rate 90\%) среднее время доступа к чанку снижается примерно в 30 раз по сравнению со сценарием без кэша. При горячем кэше (hit rate 99\%) достигается ускорение более чем в 50 раз.

\subsubsection*{Время загрузки модели и применение формулы P2P}

Полное время загрузки модели размером $S$ складывается из суммарного AMAT по всем чанкам. При $N_c = S / s$ чанках размером $s$ и параллельной загрузке с $k$ пиров~\cite{li2020caching}:

\begin{equation}
  T_{\mathrm{download}} = \frac{S}{B_{\mathrm{local}} + \sum_{i=1}^{k} B_i} + \mathrm{RTT}
  \label{eq:p2p}
\end{equation}
где $B_{\mathrm{local}}$ --- скорость чтения с локального диска, $B_i$ --- пропускная способность $i$-го пира.

Рассмотрим численный пример для модели Llama~3~8B ($S = 15$~ГБ):
\begin{itemize}[leftmargin=*]
  \item \textbf{Baseline}: $k = 0$, $B_{\mathrm{local}} = 0$ (диск не используется), $B_{\mathrm{HF}} = 200$~МБ/с. $T_{\mathrm{download}} = 15000 / 200 \approx 75$~с.
  \item \textbf{Peer cache, $k = 1$}: $B_{\mathrm{local}} = 1250$~МБ/с (10~Гбит/с). $T_{\mathrm{download}} = 15000 / 1250 \approx 12$~с --- ускорение в $\approx 6.3$ раза.
  \item \textbf{Peer cache, $k = 3$}: суммарная пропускная способность $3 \times 1250 = 3750$~МБ/с, но ограничена входящим каналом ноды 1250~МБ/с. $T_{\mathrm{download}} \approx 12$~с.
  \item \textbf{Local cache}: $B_{\mathrm{local}} = 7000$~МБ/с (NVMe). $T_{\mathrm{download}} = 15000 / 7000 \approx 2.1$~с --- ускорение в $\approx 36$ раз.
\end{itemize}

Таким образом, попадание в локальный кэш даёт наибольший выигрыш. При отсутствии локального кэша один пир обеспечивает ускорение в 6 раз по сравнению с загрузкой из внешнего репозитория.

\subsubsection*{Закон Амдала и предел параллельного ускорения}

Параллельная загрузка чанков с нескольких пиров аналогична параллельному выполнению задачи. Закон Амдала~\cite{amdahl_law} ограничивает максимальное ускорение $S(k)$:
\[
  S(k) = \frac{1}{\alpha + \frac{1-\alpha}{k}}
\]
где $\alpha$ --- доля последовательной части (Redis lookup, инициализация соединений), $k$ --- число параллельных пиров. При $\alpha = 0.05$ (5\% последовательной части) и $k = 5$ пирах: $S(5) = 1 / (0.05 + 0.19) \approx 4.2$. Таким образом, привлечение более 5~пиров даёт убывающую отдачу; практически оптимальным является диапазон $k = 2$--$4$ при балансе между параллелизмом и overhead на координацию.

\subsubsection*{Применение закона Литтла}

Закон Литтла~\cite{little1961}:
\begin{equation}
  L = \lambda \cdot W
  \label{eq:little-evaluation}
\end{equation}
позволяет оценить число одновременно ожидающих запросов. При интенсивности входящих запросов $\lambda = 5$~запр./с:
\begin{itemize}[leftmargin=*]
  \item Baseline ($W = 75$~с): $L = 5 \times 75 = 375$ ожидающих запросов.
  \item P2P/FUSE после прогрева ($W = 2.1$~с): $L = 5 \times 2.1 \approx 11$ ожидающих запросов.
\end{itemize}

Снижение числа ожидающих запросов в $\approx 34$ раза непосредственно транслируется в снижение нагрузки на очередь и улучшение соблюдения SLA по задержке ответа.

\subsection{Ожидаемые результаты сравнения}

\begin{table}[htbp]
\centering
\caption{Сравнение подходов по ключевым показателям}
\label{tab:comparison}
\small
\begin{tabularx}{\textwidth}{lXXX}
\toprule
\textbf{Показатель} & \textbf{emptyDir} & \textbf{NFS} & \textbf{P2P/FUSE} \\
\midrule
$T_{\mathrm{model\_ready}}$, первый запуск & $\approx$75~с & $\approx$75~с & $\approx$75~с (fallback) \\
$T_{\mathrm{model\_ready}}$, повтор на той же ноде & $\approx$75~с & $\approx$12~с & $\approx$2~с (local cache) \\
$T_{\mathrm{model\_ready}}$, повтор на соседней ноде & $\approx$75~с & $\approx$12~с & $\approx$12~с (peer cache) \\
$T_{\mathrm{first\_byte}}$ & $\approx$75~с & $\approx$75~с & $<$1~с (lazy FUSE) \\
Внешний трафик, $N=10$ реплик & $10 \times 15 = 150$~ГБ & $15$~ГБ & $\leq 15$~ГБ \\
SPOF & Нет & NFS-сервер & Redis (с HA~---~нет) \\
Деградация при отказе & Нет & Полная & Graceful (local+fallback) \\
Сложность эксплуатации & Низкая & Средняя & Средняя \\
\bottomrule
\end{tabularx}
\end{table}

Ключевым преимуществом P2P/FUSE-подхода является значение $T_{\mathrm{first\_byte}} < 1$~с: FUSE-монтирование позволяет runtime-контейнеру начать чтение сразу после запуска, тогда как физическая загрузка чанков происходит параллельно с чтением (lazy loading). Для baseline и NFS время до первого байта совпадает с полным временем загрузки, поскольку \texttt{storage-initializer} должен завершить скачивание до старта основного контейнера.

По показателю внешнего трафика P2P/FUSE при прогретом кластере эквивалентен NFS: загрузка происходит один раз. При этом данные распределяются по дискам всех нод, создавая избыточность без выделенного NFS-сервера.

\begin{figure}[htbp]
\centering
\includegraphics[width=0.88\textwidth]{chart_tmodel_ready.png}
\caption{Время готовности модели ($T_{\mathrm{model\_ready}}$) по сценариям для Llama~3~8B (15~ГБ). При наличии локального кэша P2P/FUSE-подход обеспечивает $T_{\mathrm{model\_ready}} = 1.9$~с против 74.8~с у baseline.}
\label{fig:tmodel_ready}
\end{figure}

\begin{figure}[htbp]
\centering
\includegraphics[width=0.82\textwidth]{chart_external_traffic.png}
\caption{Суммарный внешний трафик при масштабировании до $N$ реплик (Llama~3~8B, 15~ГБ). Подходы NFS и P2P/FUSE загружают модель однократно независимо от числа реплик; emptyDir генерирует $N \times 15$~ГБ трафика.}
\label{fig:external_traffic}
\end{figure}

\subsection{Анализ сценариев отказов}

\subsubsection*{Падение worker-ноды с чанками}

При отказе ноды агент перестаёт обновлять TTL-записи в Redis. По истечении TTL (типично 30~с) записи удаляются. При запросе затронутого чанка агент соседней ноды получает пустой список пиров для данного чанка и инициирует fallback на HuggingFace. Деградация ограничена временем $T_{\mathrm{TTL}}$ и затрагивает только чанки, хранившиеся исключительно на упавшей ноде. При репликации каждого чанка на $r \geq 2$ ноды вероятность потери доступности:
\[
  P_{\mathrm{unavail}} = p^r
\]
где $p$ --- вероятность отказа одной ноды. При $p = 0.01$ и $r = 2$: $P_{\mathrm{unavail}} = 10^{-4}$; при $r = 3$: $P_{\mathrm{unavail}} = 10^{-6}$~\cite{ford2010availability}.

\subsubsection*{Устаревшая запись в Redis (stale metadata)}

Агент пытается подключиться к указанному пиру, получает ошибку соединения, переходит к следующему пиру из множества \texttt{chunk:peers:\{hash\}} или инициирует fallback. Stale-записи автоматически удаляются по TTL, не требуя ручного вмешательства.

\subsubsection*{Недоступность Redis}

При недоступности Redis агент теряет возможность находить пиров. Система деградирует до двухуровневой: local cache~$\to$~remote fallback. Обслуживание запросов продолжается при наличии локальных данных. При развёртывании Redis Sentinel (1 master, 2 replica) вероятность одновременного отказа всех узлов пренебрежимо мала при независимых отказах~\cite{redis_sentinel_docs}.

\subsubsection*{Недоступность HuggingFace Hub}

Система использует данные из local и peer cache. Новые модели, ещё не загруженные в кластер, становятся недоступны. Это является ожидаемой деградацией при условии предварительного прогрева популярных моделей. Возможным улучшением является поддержка зеркальных репозиториев в качестве альтернативных источников fallback.

\subsubsection*{Частичная запись чанка (crash во время записи)}

При аварийном завершении агента в процессе записи чанка временный файл (с суффиксом \texttt{.tmp}) не переименовывается в финальный путь. При следующем запросе агент не найдёт чанк в локальном хранилище и повторит загрузку. Консистентность обеспечивается атомарностью системного вызова \texttt{os.Rename}, который переименовывает файл только при успешном завершении записи.

\subsubsection*{Сетевое разбиение кластера}

При разбиении кластера на изолированные части агенты, не имеющие доступа к Redis master, переходят в degraded-режим: local cache + remote fallback. После восстановления сетевой связности агенты возобновляют heartbeat, и Redis автоматически обновляет реестр активных пиров.

\begin{figure}[htbp]
\centering
\includegraphics[width=0.78\textwidth]{chart_reliability.png}
\caption{Деградация инференса при трёх сценариях отказов (прогретый кластер, $r = 3$). NFS полностью теряет доступ при отказе сервера или HF Hub; P2P/FUSE деградирует до 15\% при недоступности Redis и остаётся полностью работоспособен при сбое HF Hub.}
\label{fig:reliability}
\end{figure}

\subsection{Ограничения предложенного подхода}

\textbf{Накладные расходы FUSE.} Каждая операция чтения проходит через два context switch: пользовательское пространство~$\to$~ядро (VFS-перехват)~$\to$~пользовательское пространство (FUSE-демон). Для ML-workload с последовательными крупными чтениями (размер чанка 16~МБ) эти накладные расходы незначительны: при 2 context switch по $\approx$5~мкс каждый и времени чтения чанка $\approx$2.3~мс суммарные накладные расходы не превышают 0.5\%. Для сценариев случайного доступа по малым диапазонам (например, чтение отдельных тензоров модели) относительные накладные расходы возрастают.

\textbf{Redis как критический компонент.} Без HA-конфигурации Redis является SPOF для функции peer discovery. Данные при этом не теряются: система деградирует до local + fallback, сохраняя частичную работоспособность. HA-конфигурация (Sentinel или Cluster) устраняет эту единственную точку отказа на уровне сигнализации.

\textbf{Сложность наблюдаемости.} P2P-трафик менее прозрачен, чем чтение из NFS. Для корректного мониторинга требуется сбор метрик: счётчики local hit/miss, peer hit/miss, remote fallback; задержки по уровням; состояние Redis; объём данных на каждой ноде. Без инструментирования диагностика деградации производительности затруднена.

\textbf{Первый запуск в кластере.} При полном отсутствии кэша (cold cluster) система вынуждена загружать модель из внешнего репозитория. Выигрыш появляется начиная со второго запуска. Возможным улучшением является предварительное прогревание кэша (cache warming) в периоды низкой нагрузки.

\textbf{Малые модели.} Для моделей размером до 1~ГБ накладные расходы на FUSE, Redis и P2P-координацию могут быть сравнимы с временем прямой загрузки через \texttt{emptyDir}. Архитектура наиболее эффективна для LLM и других крупных моделей ($>$5~ГБ).

\textbf{Безопасность.} В текущем прототипе TCP-соединения между агентами не аутентифицированы. В production-окружении необходима поддержка mTLS между агентами, авторизация запросов к Redis через ACL и шифрование трафика.

\subsection{Выводы по разделу}

Расчётная модель AMAT и анализ формулы P2P-загрузки демонстрируют, что предложенная архитектура наиболее эффективна в сценариях повторных запусков и burst-масштабирования. При прогретом кластере (hit rate 90\%) среднее время доступа к чанку снижается в 30 раз по сравнению с постоянным обращением к внешнему репозиторию, а полное время загрузки модели --- с 75~с до 2~с при попадании в локальный кэш. Механизм FUSE lazy loading обеспечивает время до первого байта менее 1~с, что принципиально недостижимо при \texttt{storage-initializer}-подходе. Анализ сценариев отказов показывает, что система деградирует gracefully: потеря worker-ноды затрагивает только нерезервированные чанки, а недоступность Redis переводит систему в двухуровневый режим с сохранением работоспособности. Ключевые ограничения --- FUSE overhead при мелком доступе, требование HA Redis и сложность наблюдаемости --- известны и могут быть устранены на следующих этапах разработки.
